% External-FDT: Cloud-to-Building Control System Tunneling in BACnet Networks
% Empirical analysis of internet-exposed building automation infrastructure

\documentclass[11pt,a4paper]{article}

% Core packages
\usepackage[utf8]{inputenc}
\usepackage[T1]{fontenc}
\usepackage{lmodern}
\usepackage{microtype}
\usepackage{amsmath,amssymb,amsfonts}
\usepackage[numbers,sort&compress]{natbib}
\bibliographystyle{unsrtnat}
\usepackage{hyperref}
\usepackage{url}
\usepackage{booktabs}
\usepackage{graphicx}
\usepackage[margin=1in]{geometry}
\usepackage{xcolor}

\hypersetup{
  pdfauthor={Cameron Warren},
  pdftitle={External-FDT: Cloud-to-Building Control System Tunneling},
  colorlinks=true,
  linkcolor=blue,
  citecolor=blue
}

% Title
\title{\textbf{External-FDT}: Cloud-to-Building Control System Tunneling in BACnet Networks}
\author{Cameron Warren\\
  \texttt{cwarre33@uncc.edu}\\[0.5em]
  Department of Computer Science\\
  University of North Carolina at Charlotte
}
\date{\today}

\begin{document}

\maketitle

\begin{abstract}
Building automation systems increasingly bridge private facility networks to public 
cloud infrastructure, creating undocumented attack surfaces. Through longitudinal 
analysis of internet-exposed BACnet Broadcast Management Devices (BBMDs), we 
discover a systemic pattern we term \textbf{External-FDT}: 71\% of publicly visible 
BBMDs maintain Foreign Device Table entries pointing to public, routable IP addresses. 
Of these, 42\% tunnel to major cloud providers (AWS, Azure, DigitalOcean). We also 
discover a novel \textbf{integrator concentration} phenomenon where a single external 
endpoint serves multiple distinct building networks. Our empirical analysis of 17 seeded 
BBMDs over 47 days exposes potential supply-chain attack surfaces in building 
automation infrastructure. All findings derive from passive OSINT analysis.
\end{abstract}

\section{Introduction}
\label{sec:intro}

Building automation systems---responsible for HVAC, access control, and 
critical facility operations---have historically relied on isolated networks. 
The BACnet protocol enables these systems to span subnets through Broadcast 
Management Devices (BBMDs) that forward traffic to registered Foreign Devices.

We discover that the majority of internet-exposed BBMDs maintain Foreign 
Device registrations with \textbf{public, routable IP addresses}---creating 
undocumented bridges to external infrastructure including major cloud providers.

\subsection{Contributions}

This paper makes three empirical contributions:

\begin{enumerate}
  \item \textbf{External-FDT characterization}: We quantify the prevalence (71\%) 
    and patterns of public-IP Foreign Device registrations.
    
  \item \textbf{Integrator Concentration}: We identify a novel pattern where a 
    single external endpoint serves multiple distinct building networks.
    
  \item \textbf{Attribution methodology}: We validate passive OSINT methods 
    for identifying cloud-provider infrastructure without active probing.
\end{enumerate}

\section{Background}
\label{sec:background}

BACnet BBMDs enable cross-subnet communication by forwarding broadcasts to 
registered Foreign Devices. A Foreign Device Table entry contains: IP address, 
port, and TTL.

\subsection{Related Work}

Prior ICS research focused on protocol vulnerabilities and port exposure. No 
prior work has characterized tunneling relationships between building control 
systems and external infrastructure.

\section{Methodology}
\label{sec:methodology}

\textbf{Ethical constraints}: All analysis uses passive OSINT only---Shodan-indexed 
banners. No active connection attempts.

\textbf{Data}: Longitudinal Shodan dataset: 17 seeded BBMDs, March--April 2026 (47 days).

\textbf{Classification}: We categorize FD IPs as Internal (RFC1918) or External 
(public routable), with External further classified by cloud provider attribution.

\section{Results}
\label{sec:results}

\subsection{Prevalence}

Of 17 analyzed BBMDs, 12 (71\%) exhibited External-FDT behavior.

\subsection{Cloud Distribution}

\begin{table}[h]
\centering
\begin{tabular}{lc}
\toprule
\textbf{Provider} & \textbf{Count} \\
\midrule
DigitalOcean & 4 \\
Amazon AWS & 2 \\
Microsoft Azure & 1 \\
Other / Unknown & 5 \\
\bottomrule
\end{tabular}
\caption{Cloud Provider Attribution}
\label{tab:clouds}
\end{table}

DigitalOcean shows unexpected dominance (33\%).

\subsection{Integrator Concentration}

Two distinct BBMDs share a common Foreign Device:

\begin{table}[h]
\centering
\begin{tabular}{ll}
\toprule
\textbf{Attribute} & \textbf{Value} \\
\midrule
BBMD 1 & 66.58.248.125 \\
BBMD 2 & 24.237.132.230 \\
Shared FD & 216.67.73.166 \\
Separation & ~4,000 km \\
Duration & 47 days continuous \\
\bottomrule
\end{tabular}
\caption{Integrator Concentration Pattern}
\end{table}

This suggests a single integrator monitoring station bridging multiple 
client facility networks---creating a supply-chain attack surface.

\section{Implications}
\label{sec:implications}

External-FDT creates three attack surfaces:
\begin{enumerate}
  \item Cloud compromise yields building network access
  \item Single integrator breach affects multiple facilities
  \item Internet-exposed FD endpoints compound risk
\end{enumerate}

\section{Conclusion}

We demonstrate that cloud-to-building BACnet tunneling is systemic. The 
integrator concentration pattern suggests novel supply-chain attack surfaces 
in critical infrastructure.

\section*{Acknowledgments}

This research used only publicly available Shodan data. All findings 
reported to US-CERT for coordinated disclosure.

\end{document}
